
\medskip

%Un jeu vidéo récompense par un objet tiré au sort les joueurs ayant remporté un défi. L'objet tiré peut être \og commun \fg{} ou \og rare \fg. Deux types d'objets communs ou rares sont disponibles, des épées et des boucliers.
%
%Les concepteurs du jeu vidéo ont prévu que :
%
%\begin{itemize}
%\item la probabilité de tirer un objet rare est de 7\,\% ;
%\item si on tire un objet rare, la probabilité que ce soit une épée est de 80\,\% ;
%\item si on tire un objet commun, la probabilité que ce soit une épée est de 40\,\%.
%\end{itemize}
%
%\emph{Les parties \textbf{\rm A} et \textbf{\rm B} sont indépendantes.}
%
%\bigskip

\textbf{Partie A}

\medskip

%Un joueur vient de remporter un défi et tire au sort un objet. On note :
%
%\begin{itemize}[label=$\bullet~$]
%\item $R$ l'évènement \og le joueur tire un objet rare\fg{} ;
%\item $E$ l'évènement \og le joueur tire une épée \fg{} ;
%\item $\overline{R}$ et $\overline{E}$ les évènements contraires des évènements $R$ et $E$.
%\end{itemize}
%
%\medskip

\begin{enumerate}
\item %Dresser un arbre pondéré modélisant la situation, puis calculer $P(R \cap E)$.
On dresse l'arbre pondéré de probabilités en utilisant les données de l'énoncé :

$p(R) = 0,07 \:; p_R(E) = 0,8$ et $p_{\overline{R}}(E) = 0,4$:

\begin{center}
\pstree[treemode=R,levelsep=3cm]{\TR{}}
{\pstree{\TR{$R~~$}\taput{0,07}}
	{\TR{$E$}\taput{0,8}
	\TR{$\overline{E}$}\tbput{\red 0,2}
	}
\pstree{\TR{$\overline{R}~~$}\tbput{\red 0,93}}
	{\TR{$E$}\taput{0,4}
	\TR{$\overline{E}$}\tbput{\red 0,6}
	}
}
\end{center}

On a $p(R \cap E) = p(R) \times p_{R}(E) = 0,07 \times 0,8 = 0,056$.
\item %Calculer la probabilité de tirer une épée.
D'après la loi des probabilités totales :

$p(E) = p(R \cap E) + p\left(\overline{R} \cap E\right)$.

Or $p\left(\overline{R} \cap E\right) = p\left(\overline{R}\right) \times p_{\overline{R}}(E) = 0,93 \times 0,4 = 0,372$.

Donc $p(E) =0,056 + 0,372 = 0,428$.
\item %Le joueur a tiré une épée. Déterminer la probabilité que ce soit un objet rare. Arrondir le résultat au millième.
On a $p_{E}(R) = \dfrac{p(E \cap R)}{p(E)} = \dfrac{p(R \cap E)}{p(E)} = \dfrac{0,056}{0,428} = \dfrac{56}{428} = \dfrac{14}{107} \approx \np{0,1308}$ soit 0,131 au millième près.
\end{enumerate}

\bigskip

\textbf{Partie B}

\medskip

%Un joueur remporte $30$ défis.
%
%On note $X$ la variable aléatoire correspondant au nombre d'objets rares que le joueur obtient après avoir remporté $30$ défis. Les tirages successifs sont considérés comme indépendants.

\medskip

\begin{enumerate}
\item %Déterminer, en justifiant, la loi de probabilité suivie par la variable aléatoire $X$. Préciser ses paramètres, ainsi que son espérance.
Les évènements étant indépendants et la probabilité d'obtenir un objet rare étant toujours égale à 0,07, la variable aléatoire $X$ suit une loi binomiale $\mathcal{B}(n,\:p)$ avec $n = 30$ et $p = 0,07$.

L'espérance mathématique est $E = n \times p = 30 \times 0,07 = 2,1$.
\item %Déterminer $P(X < 6)$. Arrondir le résultat au millième.

La calculatrice donne $P(X < 6) \approx \np{0,9838}$ soit 0,984 au millième près.
\item %Déterminer la plus grande valeur de $k$ telle que $P(X \geqslant k) \geqslant 0,5$. Interpréter le résultat dans le contexte de l'exercice.
Quel que soit $k \in \N$, \: $p(X \geqslant k + 1) = 1 - p(X \leqslant k)$.
D'après la calculatrice :

$p(X \geqslant 1 + 1) = 1 - p(X \leqslant 1) \approx 0,631$, et $p(X \geqslant 2 + 1) = 1 - p(X \leqslant 2) \approx 0,351$, on a donc $k = 2$.

Dans le cadre du jeu ceci signifie que la probabilité d'obtenir au moins 2 objets rares est supérieure à $\dfrac12$.
\item %Les développeurs du jeu vidéo veulent proposer aux joueurs d'acheter un \og ticket
%d'or \fg{} qui permet de tirer $N$ objets. La probabilité de tirer un objet rare reste de 7\,\%.

%Les développeurs aimeraient qu'en achetant un ticket d'or, la probabilité qu'un joueur obtienne au moins un objet rare lors de ces $N$ tirages soit supérieure ou égale à $0,95$.
%
%Déterminer le nombre minimum d'objets à tirer pour atteindre cet objectif. On veillera à détailler la démarche mise en œuvre.
Il faut donc trouver $X$ tel que $p(X \geqslant 1) \geqslant 0,95$ ou encore $1 - p(X = 0) \geqslant 0,95 \iff p(X = 0) \leqslant 1 - 0,95 \iff p(X = 0) \leqslant 0,05$.

Or $p(X = 0) = \binom{N}{0} \times 0,07^0 \times (1 - 0,07)^N = 0,93^N$.

Il faut donc résoudre l'inéquation :

$0,93^N \leqslant 0,05 \Rightarrow N \ln 0,93 \leqslant \ln 0,05$ par croissance de la fonction logarithme népérien  et enfin $N \geqslant \dfrac{\ln 0,05}{\ln 0,93}$, car $\ln 0,93 < 0$ et son inverse $\dfrac{1}{\ln 0,93}$ aussi.

D'après la calculatrice $\dfrac{\ln 0,05}{\ln 0,93} \approx 41,3$.

Conclusion Il faut que $N \geqslant 42$.
\end{enumerate}


